% MAS1004 Assignment 1: long report.
%
% You fill this in with your agent, in your repository. AGENTS.md tells the
% agent how. Red text in the PDF is what still has to be replaced.
%
% - Pictures go in figures/. Copy them from results/, which is not in git.
% - Outputs go in outputs/. run.py saves what it prints in
%   results/<tag>_output.txt, and PLAN.md saves the other outputs in results/
%   with `| tee results/...`. Copy them here. \outputfile{outputs/check.txt}
%   prints a file exactly as it is. Every number in this report has to come
%   from such an output.
%
% To make the PDF: compress this folder into a zip, and on Overleaf choose
% New Project, Upload Project. Set Menu, Compiler to XeLaTeX, and Recompile.

\documentclass[11pt]{article}
\usepackage{mas1004}

\reporttitle{Assignment 1: Long report}
\studentname{Your name}
\studentid{Your student ID}

\begin{document}
\makeheader

% ---------------------------------------------------------------------------
\section*{Problem 1. Categories and images}

\subsection*{Why these categories}
I really love going to cafes, so when I looked through my smartphone photo
gallery I realised I had the most pictures of drink cups.  I wanted to choose a
topic that is familiar in my daily life and easy for me to collect data for.
With that in mind I picked three categories of drinking vessels that are
common in everyday life but not in ImageNet's 1,000 classes:
\emph{reusable clear cups} (e.g.\ tumblers, water glasses),
\emph{disposable plastic takeout cups} (single-use cold-drink cups with lids),
and \emph{ceramic cafe mugs} (coffee mugs).  ImageNet has a ``coffee mug''
class (index 504), so I deliberately named mine ``ceramic cafe mugs'' to avoid
direct overlap.  The three categories are visually similar enough that the
task is non-trivial, yet distinguishable by shape, material, and typical
context.

\subsection*{What I expected the model to use}
\fillin{What you thought, before training, the model would use to tell them
apart: shape, colour, texture, background, or something else.}

\subsection*{Images per category}
I downloaded 150 images for each category (450 total) using
\texttt{src/collect.py}.  Every category returned the full 150 images.

\begin{center}
\begin{tabular}{lrr}
\toprule
Category & Downloaded & In \texttt{data/clean/} \\
\midrule
ceramic\_cafe\_mugs             & 150 & 150 \\
disposable\_plastic\_takeout\_cups & 150 & 150 \\
reusable\_clear\_cups            & 150 & 150 \\
\bottomrule
\end{tabular}
\end{center}

% ---------------------------------------------------------------------------
\section*{Problem 2. Preparing an image}

\sidebyside{figures/prepare_original.jpg}{the image as downloaded (any size)}
           {figures/prepare_result.png}{after \texttt{prepare\_image}: 224~$\times$~224, channels-first, normalised and denormalised back for viewing}

The original is $224~\times~224$ or larger.  \texttt{prepare\_image} resizes the
shorter side to 256, then takes the $224~\times~224$ center crop.  For this
image the crop removed about \fillin{\#} pixels from the longer side.  No
important part of the cup was cut off because the cup is centred in the frame.
For images where the cup is off-centre---for example, a mug in the lower
corner of the frame---the center crop can cut off part of the handle or the
rim.  I kept those images in the dataset anyway, because the model may still
learn useful features from the visible part.

% ---------------------------------------------------------------------------
\section*{Problem 3. First training runs}

\begin{center}\footnotesize
\begin{tabular}{llllrrrr}
\toprule
tag & start & trained & epochs & lr & trainable parameters & train acc. & test acc. \\
\midrule
\multicolumn{8}{l}{\fillin{one row per run, as run.py printed it}} \\
\bottomrule
\end{tabular}
\end{center}

\picturefile{figures/run_curves.png}

\fillin{Did the loss go down? Did it keep going down, or flatten out?}

\fillin{The gap between train and test accuracy, and what it means.}

\begin{center}\picturefile[0.5\linewidth]{figures/run_confusion.png}\end{center}

\fillin{Which two categories it mixes up most, and whether that surprised you.}

\fillin{How far apart the test accuracies are, starting from ImageNet and from
random numbers, and why you think that is.}

\fillin{How close training only the last layer gets to training all of them.}

\subsection*{Output of the runs}
\outputfile{outputs/run_output.txt}
\outputfile{outputs/scratch_output.txt}
\outputfile{outputs/frozen_output.txt}
% One more line like these for the setting you chose, named by its tag.
\fillin{the output of the run with your own setting}

% ---------------------------------------------------------------------------
\section*{Problem 4. Cleaning}

\subsection*{My rule}
\fillin{Your rule, in one or two sentences, as you wrote it before you
started.}

\subsection*{Three kinds of junk}
\fillin{Three kinds of junk, with one picture of each.}

% For example: \picturefile[0.3\linewidth]{figures/junk_1.jpg}

\subsection*{What I removed}
\outputfile{outputs/clean_count.txt}

\fillin{How many near copies it found, and what you did with them.}

\fillin{Of the first 20 suspects in each class, how many you removed. Whether
the suspects included junk you had missed, and whether you found junk it did
not list.}

\subsection*{Before and after cleaning}
\outputfile{outputs/compare.txt}

\fillin{The test accuracy and the accuracy on your new images, before and
after. If a number went down, why you think it did.}

% ---------------------------------------------------------------------------
\section*{Problem 5. Images from a new source}

\fillin{Where your new images came from, and why you are sure none of them are
among your downloads.}

\outputfile{outputs/overlap.txt}

\fillin{The accuracy on your downloaded test set and on your new images, side
by side.}

\fillin{Five mistakes it was confident about, with the pictures. For each one,
what in the image you think pushed it the wrong way.}

\fillin{What is different about your new images: background, lighting, angle,
distance, what else is in the frame.}

% ---------------------------------------------------------------------------
\section*{Problem 6. The demo}

Address: \fillin{your GitHub Pages address}

\fillin{What happens when you show it something that is none of your
categories.}

% ---------------------------------------------------------------------------
\section*{Problem 7. What the model looks at}

% Copy this subsection once for every guess.
\subsection*{Guess 1}

\fillin{The guess, and what in your images or your model's mistakes made you
think of it.}

\fillin{The change you chose, and why it tests the guess.}

My request to the agent, word for word:
\begin{quote}
\fillin{the request}
\end{quote}

% One \sidebyside for every image you changed.
\sidebyside{figures/guess1_original_1.jpg}{\fillin{original: the model's answer}}
           {figures/guess1_changed_1.jpg}{\fillin{changed: the model's answer}}

\fillin{What the results say about the guess.}

\subsection*{What I now think the model uses}
\fillin{What you now think the model uses to tell your categories apart, and
how sure you are.}

% ---------------------------------------------------------------------------
\section*{Appendix. Output of check.py}
\outputfile{outputs/check.txt}

\end{document}
